\documentclass[aspectratio=169]{beamer}
\usetheme{Madrid}
\usecolortheme{default}
\usefonttheme{professionalfonts}
\setbeamertemplate{navigation symbols}{}
\setbeamertemplate{footline}[frame number]

\usepackage[utf8]{inputenc}
\usepackage[T1]{fontenc}
\usepackage[english]{babel}
\usepackage{amsmath,amssymb,booktabs,array}
\usepackage{graphicx}
\usepackage{tikz}
\usepackage{xcolor}
%\usepackage{siunitx}
\usepackage{ragged2e}

\title{Modelim Dinamik i Trajtimit të Ujërave të Ndotura Spitalore}
\subtitle{Nga një prototip MATLAB në një paketë Python më realiste dhe më e kontrollueshme nga terminali}
\author{Klaudio Peqini}
\institute{Material prezantues për diskutim teknik me profesorin}
\date{\today}

\newcommand{\code}[1]{\texttt{#1}}

\newcommand{\safeinclude}[2][0.72\textwidth]{%
  \IfFileExists{#2}{\includegraphics[width=#1]{#2}}{%
    \fbox{\parbox[c][0.48\textheight][c]{0.84\textwidth}{\centering
    Figura \texttt{#2} nuk u gjet gjatë kompilimit.\\[0.5em]
    Vendoseni figurën në të njëjtin direktor ose përditësoni path-in në skedarin \texttt{.tex}.}}}%
}

\begin{document}

\begin{frame}
  \titlepage
\end{frame}

\begin{frame}{Motivimi shkencor dhe inxhinierik}
\justifying
Ujërat e ndotura spitalore përmbajnë një përzierje komplekse të:
\begin{itemize}
    \item lëndës organike biodegraduese dhe jo-biodegraduese,
    \item lëndëve ushqyese (azot, fosfor),
    \item mikrondotësve farmaceutikë,
    \item dezinfektantëve dhe komponimeve inhibuese,
    \item grimcave pezull, vajrave dhe komponimeve të adsorbueshme.
\end{itemize}
Qëllimi i modelit nuk është vetëm të riprodhojë prirje numerike, por të shërbejë si \alert{mjet interpretues} për:
\begin{itemize}
    \item vlerësimin e performancës së trenit të trajtimit,
    \item analizën e skenarëve të ngarkesës,
    \item identifikimin e pikave ku sistemi dështon ose humbet efikasitetin.
\end{itemize}
\end{frame}

\begin{frame}{Objektivi i versionit të ri të paketës}
\justifying
Versioni i ri në Python u ndërtua për të adresuar kufizimet e prototipit paraprak dhe për të ofruar:
\begin{enumerate}
    \item \textbf{dinamikë më realiste të influentit},
    \item \textbf{më shumë variabla dinamikë},
    \item \textbf{trajtim të diferencuar sipas proceseve},
    \item \textbf{kontroll nga terminali} përmes \code{argparse},
    \item \textbf{figura dhe tabela} gati për raportim dhe diskutim.
\end{enumerate}
Arkitektura e trajtimit të simuluar është:
\[
\mathrm{EQ} \rightarrow \mathrm{MBR} \rightarrow \mathrm{AOP} \rightarrow \mathrm{GAC}
\]
ku:
\begin{itemize}
    \item EQ = rezervuar i barazimit,
    \item MBR = reaktor biologjik me membranë,
    \item AOP = proces oksidimi i avancuar,
    \item GAC = adsorbim në karbon aktiv granular.
\end{itemize}
\end{frame}

\begin{frame}{Ndryshimet kryesore kundrejt modelit fillestar}
\small
\begin{tabular}{>{\RaggedRight\arraybackslash}p{0.31\textwidth} >{\RaggedRight\arraybackslash}p{0.29\textwidth} >{\RaggedRight\arraybackslash}p{0.31\textwidth}}
\toprule
Aspekti & Prototipi fillestar & Paketa e re në Python \\
\midrule
Influenti & Profil sinusoidal idealizuar & Profil ditor spitalor + zhurmë + ngjarje ``shock'' \\
Blloku EQ & Vetëm konceptual & Dinamik i modeluar në mënyrë eksplicite \\
COD & Një variabël e vetme & I ndarë në fraksion biodegradues dhe inert \\
Nutrientët & Trajtim minimal & NH$_4^+$, NO$_3^-$, PO$_4^{3-}$ me transformime të thjeshtuara \\
Komponimet inhibuese & Pothuajse të munguara & Fenole, formaldehid, glutaraldehid, detergjentë, vajra, Pb \\
AOP/GAC & Shumë të përmbledhur & Parametra më specifikë sipas llojit të komponimit \\
Ekzekutimi & Skript statik & CLI nga terminali me argumente \\
\bottomrule
\end{tabular}
\end{frame}

\begin{frame}{Variablat e modeluara në paketën e re}
\justifying
Gjendja e sistemit u zgjerua për të përfshirë treguesit kryesorë të cilësisë së ujit dhe komponimeve problematike:
\begin{columns}[T]
\column{0.48\textwidth}
\begin{itemize}
    \item BOD
    \item COD biodegradues
    \item COD inert
    \item TOC
    \item NH$_4^+$
    \item NO$_3^-$
    \item PO$_4^{3-}$
    \item SS
\end{itemize}
\column{0.48\textwidth}
\begin{itemize}
    \item alkalinitet
    \item përçueshmëri
    \item fenole
    \item formaldehid
    \item glutaraldehid
    \item detergjentë anionikë
    \item vajra/yndyrna
    \item Pb, CBZ, DCF
\end{itemize}
\end{columns}
\vspace{0.5em}
Kjo e zhvendos modelin nga një qasje ``toy model'' drejt një \alert{modeli shumë-komponentësh}, më afër realitetit të ujërave të ndotura spitalore.
\end{frame}

\begin{frame}{Bërthama matematikore e modelit}
\justifying
Për rezervuarin e barazimit dhe reaktorin biologjik përdoret bilanci dinamik i masës:
\[
\frac{dC_i}{dt}=\frac{Q}{V}\left(C_{i,\,in}-C_i\right)-r_i(C,\,pH,\,I,\,F)
\]
ku:
\begin{itemize}
    \item $C_i$ është përqendrimi i komponentit $i$,
    \item $Q/V$ kontrollon kohën e qëndrimit,
    \item $r_i$ përfaqëson heqjen/transformimin biologjik,
    \item $I$ përmbledh efektet inhibuese,
    \item $F$ përfaqëson indeksin empirik të fouling-ut.
\end{itemize}
Për AOP përdoret një zhdukje efektive e rendit të parë:
\[
C_{i,out}=C_{i,in}\exp(-k_{AOP,i}\tau)
\]
Ndërsa për GAC përdoret një formulim i tipit breakthrough, i varur nga afiniteti adsorptiv dhe koha e shërbimit të shtratit.
\end{frame}

\begin{frame}{Si u bë influenti më realist}
\justifying
Në vend të një sinusoide të vetme, influenti u ndërtua si superpozim i tre elementeve:
\[
C_{in}(t)=C_{base}(t)\,[1+\epsilon(t)] + C_{shock}(t)
\]
ku:
\begin{itemize}
    \item $C_{base}(t)$ = profil ditor i ngarkesës spitalore (natë, mëngjes, pik, pasdite),
    \item $\epsilon(t)$ = luhatje të rastësishme për të shmangur një sinjal tepër të rregullt,
    \item $C_{shock}(t)$ = ngjarje impulsive që imitojnë p.sh. shkarkime dezinfektantësh ose ngarkesa laboratorike.
\end{itemize}
\vspace{0.5em}
Kjo zgjedhje jep kurba më të besueshme fizikisht dhe shmang prodhimin e rezultateve periodike artificiale.
\end{frame}

\begin{frame}{Rezultatet kryesore që duhen diskutuar me profesorin}
\justifying
Paketa gjeneron automatikisht figura dhe tabela që ndahen në pesë grupe interpretimi:
\begin{enumerate}
    \item \textbf{Prirjet kohore të ndotësve masivë} (BOD, COD, NH$_4^+$, PO$_4^{3-}$, SS).
    \item \textbf{Prirjet e mikrondotësve} (p.sh. CBZ, DCF) përgjatë AOP dhe GAC.
    \item \textbf{Mesataret e efluentit për 24 orë} për krahasim me objektiva/toleranca.
    \item \textbf{Breakthrough i GAC} për të kuptuar humbjen graduale të kapacitetit adsorbues.
    \item \textbf{pH dhe fouling} si tregues operacionalë të qëndrueshmërisë së sistemit.
\end{enumerate}
Qëllimi interpretues nuk është një numër i vetëm, por \alert{sjellja e zinxhirit të trajtimit} nën skenarë të ndryshëm.
\end{frame}

\begin{frame}{Figura 1: Parametrat kryesorë të cilësisë së ujit}
\centering
\safeinclude{results_trends.png}
\vspace{0.4em}

\small
\justifying
Kjo figurë përdoret për të diskutuar se si komponentët bazë zbuten përgjatë trenit të trajtimit dhe si dinamika e influentit filtrohet nga EQ dhe MBR.
\end{frame}

\begin{frame}{Figura 2: Mikrondotësit dhe roli i AOP/GAC}
\centering
\safeinclude{micropollutants_trends.png}
\vspace{0.4em}

\small
\justifying
Kjo figurë ilustron pse blloqet e pas-trajtimit janë të nevojshme: reaktori biologjik nuk është i mjaftueshëm për komponimet persistente, ndërsa AOP dhe GAC japin heqjen shtesë që kërkohet.
\end{frame}

\begin{frame}{Figura 3: Efluenti mesatar 24-orësh}
\centering
\safeinclude{effluent_24h_avg.png}
\vspace{0.4em}

\small
\justifying
Mesatarja 24-orëshe është tregues shumë më i përshtatshëm për diskutim teknik sesa sinjale të çastit, sepse lejon krahasim me objektiva operacionale dhe me kufij të mundshëm të shkarkimit.
\end{frame}

\begin{frame}{Figura 4: Breakthrough në GAC dhe jetëgjatësia funksionale}
\centering
\safeinclude{gac_breakthrough.png}
\vspace{0.4em}

\small
\justifying
Kjo figurë u shërben vendimeve praktike: kur fillon humbja e efikasitetit të adsorbimit dhe sa i ndjeshëm është sistemi ndaj përmasës së shtratit, afinitetit dhe ngarkesës së hyrjes.
\end{frame}

\begin{frame}{Figura 5: pH dhe fouling si tregues operacionalë}
\centering
\safeinclude{pH_and_fouling.png}
\vspace{0.4em}

\small
\justifying
Ky panel ndihmon në diskutimin se jo çdo gjë është ``heqje përqendrimesh''; një model më realist duhet të japë edhe sinjale operacionale për stabilitetin e procesit biologjik dhe për degradimin gradual të performancës së membranës.
\end{frame}

\begin{frame}{Kontrolli nga terminali: avantazh praktik}
\justifying
Paketa e re mund të komandohet nga terminali pa modifikuar kodin burimor.
\vspace{0.4em}

Shembuj tipikë:
\begin{itemize}
    \item simulim nominal:
    \code{PYTHONPATH=src python -m hospital\_wwtp.cli simulate --scenario nominal}
    \item skenar me oksidim më të fortë:
    \code{... --oxidation-scale 1.3}
    \item skenar me ngjarje ``shock'':
    \code{... --shock-hour 42 --shock-width-h 1.0 --shock-multiplier 3.5}
    \item ekzekutim në batch për disa skenarë:
    \code{PYTHONPATH=src python -m hospital\_wwtp.cli batch}
\end{itemize}
\vspace{0.5em}
Kjo e bën paketën të dobishme jo vetëm si kod kërkimor, por edhe si mjet të shpejtë për analiza krahasuese dhe ndjeshmëri parametrash.
\end{frame}

\begin{frame}{Përfundime teknike}
\justifying
Nga analiza e modelit të ri dalin disa përfundime të rëndësishme:
\begin{itemize}
    \item Shtimi i \textbf{EQ dinamik} është thelbësor për zbutjen realiste të ngarkesave hyrëse.
    \item Ndarja e \textbf{COD në fraksione} përmirëson interpretimin e heqjes biologjike kundrejt komponimeve refraktare.
    \item \textbf{AOP dhe GAC} nuk janë blloqe luksi; ato janë të domosdoshme për mikrondotësit persistente.
    \item Komponimet inhibuese dhe fouling-u ndikojnë drejtpërdrejt në robustësinë e sistemit.
    \item Për diskutim me profesorin, figurat duhet lexuar si \alert{sjellje të sistemit nën skenarë}, jo si ``e vërtetë absolute'' pa kalibrim eksperimental.
\end{itemize}
\end{frame}

\begin{frame}{Punë e ardhshme e rekomanduar}
\justifying
Hapat më të logjikshëm për zhvillim të mëtejshëm janë:
\begin{enumerate}
    \item kalibrim me të dhëna laboratorike ose pilot,
    \item përmirësim i parametrave të AOP sipas komponimit,
    \item modelim më i kujdesshëm i adsorbimit në GAC,
    \item analizë ndjeshmërie dhe pasigurie me intervale parametrike,
    \item krahasim i skenarëve ``normal'', ``high-load'' dhe ``shock''.
\end{enumerate}
Kjo do ta transformonte paketën nga një prototip shumë i mirë inxhinierik në një mjet më të fortë për diskutim shkencor dhe dizajn paraprak.
\end{frame}

\begin{frame}{Pse paketa e mëparshme në MATLAB dha rezultate jorealiste? (I)}
\justifying
Arsyeja kryesore nuk ishte ``gabim matematikor'', por \alert{zgjedhja e input-it dhe niveli i thjeshtimit}:
\begin{itemize}
    \item influenti ishte i imponuar si \textbf{sinusoidë ideale} me periudhë të fiksuar,
    \item ekuacionet e reaktorit ishin kryesisht \textbf{lineare të rendit të parë},
    \item rezervuari EQ nuk ishte modeluar realisht si bllok dinamik,
    \item mungonin luhatjet stohastike, pikat ditore dhe ngjarjet impulsive.
\end{itemize}
Për një sistem linear të detyruar me sinjal sinusoidal,
\[
C_{in}(t)=C_0+A\sin(\omega t)
\]
është e natyrshme që zgjidhja në dalje të mbetet afërsisht sinusoidale, vetëm me amplitudë më të vogël dhe me vonesë fazore.
\end{frame}

\begin{frame}{Pse paketa e mëparshme në MATLAB dha rezultate jorealiste? (II)}
\justifying
Si pasojë, kurbat finale kishin një pamje shumë të rregullt dhe jo-tipike për ujërat e ndotura spitalore reale.
\vspace{0.4em}

\textbf{Pikat kryesore të kritikës teknike:}
\begin{itemize}
    \item ritëm periodik tepër i lëmuar,
    \item mungesë e pikave të ngarkesës dhe e episodëve ``shock'',
    \item mungesë e zbutjes reale nga EQ,
    \item model biologjik tepër i përmbledhur për të deformuar ose filtruar në mënyrë jolineare sinjalin hyrës.
\end{itemize}
\vspace{0.4em}

\textbf{Mesazhi përfundimtar:} rezultati jorealist ishte kryesisht pasojë e një \alert{prototipi paraprak me input shumë idealizuar}, jo e një ideje të gabuar modelimi. Pikërisht për këtë arsye u zhvillua versioni i ri në Python.
\end{frame}

\begin{frame}
\centering
\Huge Faleminderit!

\vspace{1em}
\Large Diskutim, komente dhe sugjerime
\end{frame}

\end{document}
